% BHU PYQ -- Manually transcribed & error-corrected
\documentclass[11pt,a4paper]{article}

\usepackage[a4paper,top=2.5cm,bottom=2.5cm,left=2.8cm,right=2.8cm,headheight=14pt]{geometry}
\usepackage{amsmath,amssymb}
\usepackage[T1]{fontenc}
\usepackage[utf8]{inputenc}
\usepackage{lmodern}
\usepackage{microtype}
\usepackage{fancyhdr}
\usepackage{enumitem}
\usepackage{hyperref}
\usepackage{lastpage}
\usepackage{parskip}

\hypersetup{colorlinks=true,urlcolor=black,linkcolor=black,
  pdftitle={CHB-361: Organic & Physical Chemistry -- BHU PYQ},pdfauthor={Banaras Hindu University}}

\pagestyle{fancy}\fancyhf{}
\renewcommand{\headrulewidth}{0.4pt}\renewcommand{\footrulewidth}{0.4pt}
\fancyhead[L]{\small   \textit{Banaras Hindu University}}
\fancyhead[R]{\small   \textit{CHB-361: Organic \& Physical Chemistry}}
\fancyfoot[C]{\small Page  \thepage\ of \pageref{LastPage}}


\newlist{parts}{enumerate}{2}
\setlist[parts,1]{label=(\alph*),leftmargin=*,itemsep=5pt,topsep=3pt}
\setlist[parts,2]{label=(\roman*),leftmargin=*,itemsep=3pt,topsep=2pt}

\newcommand{\pts}[1]{\hfill   \textit{[#1]}}


\begin{document}


\begin{center}
  {\large\bfseries BANARAS HINDU UNIVERSITY}\\[2pt]
  {
\normalsize B.Sc. (Hons.) Semester III Examination, 2016-17}\\[6pt]
  \rule{0.8\linewidth}{0.5pt}\\[6pt]
  {\large\bfseries Paper No. CHB-361: Organic \& Physical Chemistry}\\[4pt]
  {
\normalsize Time: 3 Hours\hfill Maximum Marks: 70}
\end{center}
\rule{\linewidth}{0.4pt}
\small



\noindent   \textit{Instructions: Answer five questions in total, including Question~1 from each section which is compulsory. Use separate answer books for Section~A and Section~B.}

\normalsize
\bigskip


\begin{center}
  {\large\bfseries SECTION A}\\[2pt]
  {
\normalsize (Organic Chemistry-II -- Marks: 35)}
\end{center}
\rule{0.4\linewidth}{0.2pt}
\smallskip




\noindent   \textbf{Question 1.} Answer all parts of the following: \pts{5 $ \times$ 2 = 10}

\begin{parts}
  \item Define H\"uckel's rule and show that cycloheptatrienyl cation (tropylium ion) is aromatic.
  \item Describe Gattermann aldehyde synthesis of benzaldehyde.
  \item Predict the products when benzaldehyde is reacted with: (i) Conc. $ \text{KOH}$ and (ii) Acetophenone in dilute $ \text{NaOH}$.
  \item Give the synthesis and mechanism of phenolphthalein from phthalic anhydride and phenol in the presence of $ \text{H}_2 \text{SO}_4$.
  \item Explain why pyridine is a stronger base than pyrrole.
\end{parts}

\medskip\rule{\linewidth}{0.2pt}




\noindent   \textbf{Question 2.}

\begin{parts}
  \item Outline the synthesis of Alizarin and Indigo dyes starting from phthalic anhydride. Give structural representations. \pts{6}
  \item Write short notes on: \pts{6}
  
\begin{parts}
    \item Bischler-Napieralski synthesis of isoquinoline
    \item Reimer-Tiemann reaction mechanism
  \end{parts}
\end{parts}

\medskip\rule{\linewidth}{0.2pt}




\noindent   \textbf{Question 3.}

\begin{parts}
  \item Discuss electrophilic aromatic substitution in naphthalene. Why does substitution occur preferentially at the $\alpha$-position? \pts{6}
  \item How will you synthesize: (i) Malachite green and (ii) Fluorescein dyes? \pts{6}
\end{parts}


\newpage

\begin{center}
  {\large\bfseries SECTION B}\\[2pt]
  {
\normalsize (Physical Chemistry-I -- Marks: 35)}
\end{center}
\rule{0.4\linewidth}{0.2pt}
\smallskip




\noindent   \textbf{Question 4.} Answer all parts of the following: \pts{5 $ \times$ 2 = 10}

\begin{parts}
  \item Sketch the diagram of a conductivity cell and define cell constant.
  \item State Kohlrausch's law of independent migration of ions.
  \item Explain why the equivalent conductance of a strong electrolyte increases with dilution.
  \item Give a brief definition of transport number of an ion.
  \item Define the term ionic mobility and state its units.
\end{parts}

\medskip\rule{\linewidth}{0.2pt}




\noindent   \textbf{Question 5.}

\begin{parts}
  \item Describe the moving boundary method for the determination of transport numbers of ions. \pts{6}
  \item The resistance of a $0.1\, \text{M}$ $ \text{KCl}$ solution in a cell was found to be $85\, \text{ohm}$ (specific conductivity of $0.1\, \text{M}$ $ \text{KCl} = 0.01289\, \text{ohm}^{-1}\, \text{cm}^{-1}$). Calculate the cell constant. If the resistance of a $0.05\, \text{M}$ solution of another electrolyte in the same cell is $220\, \text{ohm}$, calculate its specific and molar conductivity. \pts{6}
\end{parts}

\medskip\rule{\linewidth}{0.2pt}




\noindent   \textbf{Question 6.}

\begin{parts}
  \item Discuss the Debye-H\"uckel-Onsager equation for strong electrolytes and explain the relaxation and electrophoretic effects. \pts{8}
  \item Explain conductometric titrations of: (i) $ \text{HCl}$ vs. $ \text{NaOH}$ and (ii) $ \text{CH}_3 \text{COOH}$ vs. $ \text{NaOH}$. \pts{4}
\end{parts}

\vfill
\rule{\linewidth}{0.4pt}

\begin{center}\small
  \href{https://bhu-exam.vercel.app/}{Click here to visit our website}\\[4pt]
  \href{https://me-aryan.vercel.app/}{Click here to visit our contributor}\\[4pt]
  \href{https://quashan-me.vercel.app/}{Click here to visit our contributor}
\end{center}

\end{document}
